% Sections 19.3-19.8, from lectures/L19/appendix/b_fpga_bringup_and_review.md and
% lectures/L19/README.md.

\section{Designing \code{can_spi_node}}
\label{c19:sec:node}

\code{can_controller} is \textbf{not} the entity that goes onto the board, and it is worth being
precise about that. Its interface is at register level: \code{tx_id} is 11 bits, \code{tx_dlc} 4,
\code{tx_data} 64, plus \code{tx_req} and five status outputs. Nothing on a DE0-CV offers an
interface that wide, and in the finished system those inputs come from the driver's registers rather
than from a human hand.

\Chapref{c18:ch} built the layer making them reachable, and \secref{c19:sec:bridge} the transport in
front of it. The last module the book writes is therefore \file{can_spi_node.vhd}: the top level
tying \code{spi_slave}, \code{spi_reg_bridge}, \code{register_bank} and \code{can_controller}
together into one node, and exposing exactly two things to the outside world, SPI and the CAN bus.

Unlike every other module in the project it comes with no module-level testbench of its own, and that
is deliberate: every block inside it already has one, and what remains \textemdash{} pins, wiring and
voltage levels \textemdash{} is exactly what a simulation cannot model. The board is its testbench,
and the bring-up steps in \secref{c19:sec:bringup} are how you read it. (There is a system testbench
handed out, \code{can_spi_node_tb}, driving two complete nodes from modelled SPI masters; it is the
only one in the book that drives \code{sclk}, \code{mosi} and \code{ss} as real waveforms.)

\subsection*{Interface}

\begin{booktable}{@{}p{4.6em}p{3.4em}p{5.2em}L@{}}
  \thead{Port} & \thead{Dir.} & \thead{Type} & \thead{Meaning} \\
  \midrule
  \code{clock} & in & \code{std_logic} & The board's 50 MHz oscillator. The design's only clock. \\
  \code{reset_n} & in & \code{std_logic} & Raw, active-low reset, straight from a pin. Asynchronous;
    this level synchronises it. \\
  \code{sclk} & in & \code{std_logic} & The SPI clock from the master. Never used as a clock. \\
  \code{mosi} & in & \code{std_logic} & Out from the master, into the slave. \\
  \code{ss} & in & \code{std_logic} & Chip select, active low. Frames the transaction. \\
  \code{miso} & out & \code{std_logic} & Out from the slave, into the master. \\
  \code{rx_bus} & in & \code{std_logic} & The CAN bus's level, straight from the pin and
    asynchronous. \code{can_controller} synchronises it itself. \\
  \code{tx_bus} & out & \code{std_logic} & The bit this node wants to drive. \\
  \code{bus_en} & out & \code{std_logic} & Whether it is driving at all. \\
\end{booktable}

The three bus ports are the digital-side interface a real CAN transceiver presents
(\chapref{c10:ch}), and that is why they leave the design as three separate signals rather than one
wire. Turning them into a single physical wire is a board-level job, and \secref{c19:sec:wire} is
about that.

The order is SPI's four wires as a group, \code{sclk}/\code{mosi}/\code{ss}/\code{miso}, and then the
CAN bus's three. \code{miso} therefore sits among the SPI ports and not with the other outputs, so
that the four pins stand in the same order as in the wiring table in \secref{c19:sec:bringup}. That
matters: \code{can_spi_node_tb} binds positionally like everything else, and \code{miso} and
\code{rx_bus} have the same type, so swapping them passes analysis without a word and gives a node
that never answers.

There are no generics. Everything the design needs that a pin cannot deliver is either a constant in
\code{can_def} or a register the driver writes.

\subsection*{Behaviour: structural, and that is the point}
\code{can_spi_node} has no state machine, no protocol knowledge and almost no logic at all. It has
exactly three jobs:
\begin{itemize}
  \item \textbf{Instantiate the four blocks and wire them in a row.} \code{spi_slave} turns
    SCK/MOSI/SS into bytes; \code{spi_reg_bridge} turns bytes into register accesses;
    \code{register_bank} turns register accesses into controller-level signals;
    \code{can_controller} turns them into a frame. Every boundary is one of the contracts already
    verified by a testbench handed out, so if the chain misbehaves the fault is in a connection here,
    not in a block.
  \item \textbf{Synchronise the reset.} A single \code{meta_prev} instance takes the raw
    \code{reset_n} and produces the \code{reset_s2_n} every subblock takes. \code{can_controller} is
    the exception that proves the rule: it owns its reset synchroniser internally, alongside the one
    for \code{rx_bus}, so this level feeds it the raw \code{reset_n} and synchronises only for the
    SPI-side blocks. That makes three \code{meta_prev} instances in the whole design \textemdash{}
    two inside \code{can_controller}, one here.
  \item \textbf{Let the controller-facing ports go straight across.} \code{register_bank}'s
    \code{tx_id}, \code{tx_dlc}, \code{tx_data} and \code{tx_req} to \code{can_controller}, and its
    \code{tx_done}, \code{rx_id}, \code{rx_dlc}, \code{rx_data}, \code{rx_valid} and \code{error}
    back. Port to port, nothing in between. \Chapref{c18:ch} designed the two interfaces to mirror
    each other precisely so that this works out.
\end{itemize}

If a rule in this file ever has to mention a CAN field or an SPI byte, it belongs in another file.
That is the test to apply while you write it.

% ------------------------------------------------------------------------------------------
\section{Getting the bus out onto a wire}
\label{c19:sec:wire}

\code{tx_bus} and \code{bus_en} are a \emph{pair}, not one signal: the design says ``drive
dominant'', ``release'', never ``drive recessive''. Two nodes' push-pull outputs tied to one wire
would fight each other the moment one drove low and the other high. Two ways to make the shared wire
behave:

\paragraph{With a transceiver}
Which is what a real node uses: \code{tx_bus} goes to its TXD, its RXD comes back as \code{rx_bus},
and the transceiver does the open-drain work in the analogue domain.

\paragraph{Without one}
With two DE0-CVs wired directly together for the two-node demo, the FPGA's own I/O does the job. Pick
a GPIO pin and drive it open-drain:

\begin{vhdlcode}
GPIO_0(2) <= '0' when ((bus_en = '1') and (tx_bus = '0')) else 'Z';
rx_bus    <= GPIO_0(2);
\end{vhdlcode}

with a pull-up restoring \code{'1'} while every node has released the wire \textemdash{} the FPGA's
weak pull-up enabled on that pin in the Pin Planner, or an external resistor. That is
\secref{c10:sec:can}'s wired-AND rule made physical, in a single concurrent assignment.
\code{GPIO_0(0)} and \code{GPIO_0(1)} are worth driving with \code{tx_bus} and \code{bus_en}
unregistered as well, purely so that a logic analyser can see what the node \emph{intended} beside
what the wire actually did; every unused index is driven \code{'Z'}.

Nothing here is synchronised, for the same reason \code{reset_n} is not synchronised before it reaches
\code{can_controller}: the controller owns its own asynchronous inputs, so the pin is run straight to
the \code{rx_bus} port.

The pin indices are a convenience in the Pin Planner, not a requirement. Any three pins on the header
work, as long as two connected boards agree; see appendix~\ref{app:quartus}.

% ------------------------------------------------------------------------------------------
\section{Optional: a board-level wrapper for debugging}
\label{c19:sec:wrapper}

All of the above is reachable only over SPI, which means that when the link does not work you have no
way of seeing into it. A small wrapper with buttons and LEDs, standing beside \code{can_spi_node}
rather than replacing it, is worth an hour if the bring-up gets stuck \textemdash{} and the three
decisions it forces are worth understanding whether or not you build it.

The wrapper drives \code{register_bank}'s ports directly from buttons and switches and latches its
status outputs onto LEDs:

\lecfig[0.86]{lectures/L19/appendix/images/can_demo.png}{A board-level wrapper: buttons and switches
into the register bank, the status bits out onto LEDs.}{c19:fig:demo}

\begin{itemize}
  \item \textbf{A held button has to become a one-cycle pulse}, and not just for neatness. A button
    held for a few milliseconds is hundreds of thousands of clock cycles at 50~MHz. Feed that straight
    into \code{tx_req} and the controller starts a new frame the moment the previous one ends, over
    and over \textemdash{} and every pass through the start state wipes \code{error} before anything
    has read it. So: \textbf{synchronise} the raw pin with a \code{meta_prev}, \textbf{debounce} the
    synchronised level with a counter accepting a new value only once it has held for about 10~ms, and
    \textbf{then} edge-detect it. All three, in that order: edge-detecting before the debounce gives
    one pulse per bounce, and debouncing before the synchronisation feeds the counter a signal that
    was never safely sampled. This is the same conversion \code{register_bank} performs for
    \code{TX_SEND} (\chapref{c18:ch}), and it is why the bank makes the wrapper optional rather than
    necessary.
  \item \textbf{A switch bank is the judgement call.} Ten switches crossing into the clock domain are
    synchronised bit by bit, which removes the metastability but not the \emph{skew}: two switches
    flipped at the same time can land one clock apart, so \code{tx_id} could briefly read a value that
    was never on the switches. That is acceptable only because a human sets them and then leaves them
    alone. Connect the same ten bits to something that changes every cycle and a bit-by-bit
    synchroniser is entirely the wrong tool \textemdash{} that takes a handshake, or an encoding where
    only one bit changes at a time. Worth saying outright, since ``just widen the synchroniser'' is a
    habit that bites sooner or later.
  \item \textbf{Status LEDs need latching}, since \code{tx_done} and \code{rx_valid} are 20~ns pulses.
    That is the same problem \code{register_bank} solves for the driver, and solving it twice in two
    places is exactly how the two copies drift apart \textemdash{} so drive the LEDs from the bank's
    STATUS bits rather than from the controller's pulses, and the wrapper stays a wrapper.
\end{itemize}

% ------------------------------------------------------------------------------------------
\section{Bring-up}
\label{c19:sec:bringup}

The bring-up has two halves: the CAN side between two boards, and the SPI side between a
microcontroller hat and the board. Take them in that order. The CAN half needs no microcontroller at
all, so it can be run while the hat is still on the soldering bench, and it rules out a whole class of
faults before there is any SPI wire to suspect.

Once the design passes simulation in full, synthesise and program \code{can_spi_node}: New Project
Wizard, device \code{5CEBA4F23C7N}, Pin Planner, Compile, Programmer. The steps are in
appendix~\ref{app:quartus}, which is the authoritative reference for that flow, right down to the
add-on for Cyclone V devices and to getting the USB-Blaster visible to the Programmer. Note that
\code{CLOCK_50} is the one pin you may not choose freely.

\subsection*{The CAN half: two boards on one bus}
Two boards, their bus lines on \code{GPIO_0(2)} tied together into \textbf{one} wire with a common
ground, and the open-drain driving from \secref{c19:sec:wire} makes the shared wire a wired-AND.

\begin{enumerate}
  \item \textbf{One board alone.} The bus line carries nothing but the pull-up, so the node reads back
    exactly the level it drives: a loopback. \code{tx_bus} and \code{bus_en} should toggle field by
    field on a logic analyser via \code{GPIO_0(0)}/\code{GPIO_0(1)}, \code{tx_done} pulse at the end
    of EOF, and \code{error} never go high.
  \item \textbf{Two boards, one wire.} Press a button on one board and an LED should light on the
    other. That is the first time the design meets two \emph{independent} oscillators; in
    \code{can_controller_tb} the nodes shared one \code{clock}.
  \item \textbf{Arbitration for real.} Have both boards transmit at once and watch the one with the
    lower identifier win, and the loser back off.
\end{enumerate}

A loopback is also where \secref{c10:sec:bittiming}'s 70~\% sample point quietly earns its keep. The
node's own signal now takes a real round trip, out through a pad, along the wire, back in, and through
the \code{rx_bus} synchroniser's two flip-flops, at a cost of a handful of clock cycles. With
\code{TICKS_PER_BIT = 50} and \code{SAMPLE_TICK = 35}, the node reads the bus 35 ticks into a bit it
began driving at tick 0, so its own echo has long since settled and the arbitration monitoring never
misreads it. Sampling at the bit period's \emph{start} would instead read the echo of the previous
bit, and the arbitration monitoring (sent recessive, read dominant, \chapref{c17:ch}) would report a
false loss at the first recessive bit following a dominant one.

A loopback deliberately cannot show the node \emph{receiving} that frame. \code{can_controller}
latches \code{role = '1'} for the whole frame and \code{rx_shift_reg} runs only while
\code{role = '0'} (\chapref{c17:ch}), so a transmitter never decodes its own traffic; that is exactly
why the system testbench needs two nodes.

\textbf{The second board closes the gap \secref{c19:sub:gaps-tb} calls the largest.} The testbench's
two nodes share one \code{clock}: \code{resync} still phases them at every SOF, but their timers can
never \emph{drift apart} mid-frame. Two boards have two crystals, running at 50~MHz only within the
tolerance the components are specified to, so their bit periods genuinely differ, and the receiver
free-runs on its own timer for the rest of the frame after its single \code{resync} at SOF. Honest
arithmetic keeps expectations in check: at typical crystal tolerances of a few tens of ppm the drift
over a frame is a small fraction of a tick, so what that demo visibly proves is that the whole chain
works across genuinely independent clocks from an arbitrary starting phase; how much drift the design
would tolerate before the sample point wanders off the bit is the calculation in
Exercise~\ref{c12:ex:drift}.

\subsection*{The SPI half: a microcontroller hat against the board}
The wiring, per node. It is not yours to choose: it is in appendix~\ref{app:protocol}, since the
parallel class writes its driver against the same table.

\begin{booktable}{@{}p{6em}p{6em}L@{}}
  \thead{Signal} & \thead{AVR32DB28} & \thead{DE0-CV} \\
  \midrule
  \code{SCK} & \code{PC2} & \code{sclk} \\
  \code{MOSI} & \code{PC0} & \code{mosi} \\
  \code{MISO} & \code{PC1} & \code{miso} \\
  \code{SS} & \code{PC3} & \code{ss} \\
  I/O supply & \code{VDDIO2} & 3.3 V \\
  Core supply & \code{VDD} & (the hat's own 5 V) \\
  GND & \code{GND} & GND \\
\end{booktable}

No level shifter is needed, and that is a choice and not a simplification: the AVR32DB28's
\code{PORTC} is a voltage domain of its own, supplied from \code{VDDIO2} rather than \code{VDD}, so it
is put on the DE0-CV's 3.3~V while the core stays at 5~V. Every level the FPGA sees is then a level it
is specified for, and the 3.3~V the FPGA drives back is read against a 3.3 V threshold rather than a
5 V one. SPI0 is moved there with \code{PORTMUX.SPIROUTEA = ALT1}.

\begin{warning}
\textbf{Three requirements on the hat, and they have to be in place before it is etched.}
\begin{itemize}
  \item \textbf{\code{PC0}--\code{PC3} are reserved.} They are SPI0's ALT1 mux and must carry no other
    peripheral. No LED, no button, and above all no decoupling capacitor on \code{SS} \textemdash{} a
    button with 100~nF on the pin makes the \code{SS} edge unusable at 1~MHz.
  \item \textbf{Nothing else on \code{PORTC}.} The whole port sits in the \code{VDDIO2} domain at
    3.3~V, so a 5 V peripheral there would both stop working and feed 5~V into a 3.3 V domain.
  \item \textbf{\code{VDDIO2} is a track of its own.} It must not be tied to \code{VDD} on the board,
    but brought out to a pin where the DE0-CV's 3.3~V is connected, with 100~nF near the chip.
\end{itemize}

Whether the domain is separate at all is moreover controlled by \code{MVSYSCFG} in
\code{FUSE.SYSCFG1}: in single-supply mode \code{VDDIO2} is tied internally to \code{VDD},
\code{PORTC} becomes a 5 V port, and the wiring above then puts 5~V straight into inputs that are not
5 V tolerant \textemdash{} on a hat that looks exactly like a correctly built one. Check the fuse
before the first connection. It is the one fault in this chapter that can cost a board.
\end{warning}

The steps are taken in this order, and in no other. Every step rules out a source of failure, so that
whatever still does not work after step $n$ is certainly in step $n+1$:

\begin{enumerate}
  \setcounter{enumi}{3}
  \item \textbf{SPI loopback on the hat alone}, with \code{PC1} connected to \code{PC0} and the FPGA
    disconnected. Proves that the master works, and along the way that \code{PORTMUX} and
    \code{VDDIO2} are set correctly.
  \item \textbf{Register echo}: write \code{TX_ID}, read it back. Proves that the whole transaction
    chain carries.
  \item \textbf{\code{STATUS} polling}: read \code{STATUS} and see that bit 0 is set after reset.
  \item \textbf{One node transmits on command}: write a frame and trigger \code{TX_SEND}, and watch
    \code{STATUS} bit 0 go low and back \textemdash{} now triggered over SPI instead of by a button.
  \item \textbf{Two nodes, the whole chain}: one transmits on an SPI command, the other receives and
    acknowledges with \code{RX_ACK}. That is the system.
\end{enumerate}

\paragraph{What remains}
By the end of the session both halves have run on real hardware. What remains is not the wiring but
the \textbf{counterpart}: your node has answered a test program written to succeed, not a driver
written by somebody else against the specification alone. That is why the protocol specification's
timing requirements \textemdash{} the 60~ns around \code{SS}, the rule about what \code{MISO} carries
during the command byte, and the abort rule \textemdash{} are written as requirements and not as
advice: a test program that never does anything unexpected never confirms them, and a driver that does
will fail you on them.

% ------------------------------------------------------------------------------------------
\section{How production CAN controllers go further}
\label{c19:sec:gaps}

This book's controller is deliberately simplified: a teaching design rather than a certifiable one.
The list below is this book's version of it; \canref{7} sets out the same list from the protocol side,
together with the block overview and the three limitations reaching all the way up into the driver's
code, and \canref{6} describes the error handling \textemdash{} error frames, error counters,
error-passive and bus-off \textemdash{} that this design lacks entirely.

Notable gaps, worth naming explicitly:
\begin{itemize}
  \item \textbf{Standard data frames only:} the wire format is a bit-exact standard CAN data frame,
    RTR included, but there are no 29-bit extended identifiers (the receiver does not even check IDE,
    so an extended frame would be misinterpreted until its CRC check failed), and a remote frame (a
    recessive RTR) is rejected with \code{error} instead of being answered with the requested data.
  \item \textbf{No error frames and no bus-off recovery:} a real CAN node counts transmit and receive
    errors and enters progressively more restricted states (error-passive, bus-off) to protect a
    healthy bus from a broken node; this book's \code{can_controller} merely aborts the frame in
    progress locally.
  \item \textbf{No bit or form error detection:} this transmitter compares what it sent against what
    the bus read only during arbitration (\chapref{c17:ch}). A real controller monitors every bit it
    sends, flags a bit error the moment the bus disagrees outside the arbitration field and the ACK
    slot, and on reception checks that the tail's fixed fields have their prescribed values.
  \item \textbf{No retransmission on a missing acknowledgement:} a real controller retries
    automatically. This design does not even notice, since the transmitter never reads the ACK slot
    back.
  \item \textbf{No acceptance filtering and no receive buffering:} every frame on the bus lands in the
    same single set of \code{rx_*} ports, and a second frame arriving before the driver has read the
    first simply overwrites it. Production controllers offer identifier filters in hardware and a
    receive FIFO or a set of mailboxes. This is the gap the parallel class's driver feels most
    directly, and \chapref{c18:ch} shows exactly where it becomes visible at register level.
  \item \textbf{A fixed bit timing model:} this book samples at a fixed 70~\% point; production
    controllers expose configurable propagation and phase segments and a synchronisation jump width,
    tuned to bus length and node count.
  \item \textbf{An aborted reception (after lost arbitration) is simply discarded}, rather than
    seamlessly continuing to receive the frame that won; a real node goes on receiving whoever won the
    bus. This design's node waits out the rest of the frame in \code{STATE_IDLE}
    (\secref{c16:sec:roles}) and picks up from the next one.
  \item \textbf{A simplified bus-idle rule:} this design requires eleven recessive bit periods, the
    same number real CAN asks of a node joining bus traffic, before any node may leave
    \code{STATE_IDLE}. Real CAN is finer-grained: a node that has just taken part in a frame may
    transmit again after the 3-bit intermission following EOF, and a node that has been bus-off waits
    out 128 occurrences of eleven recessive bits before it may join again.
\end{itemize}

% ------------------------------------------------------------------------------------------
\section{The design, seen in the rear-view mirror}
\label{c19:sec:review-design}

Every chapter from \chapref{c10:ch} onwards added exactly one responsibility, never duplicated
anywhere else:

\begin{booktable}{@{}p{9.6em}Lp{5.4em}@{}}
  \thead{Block} & \thead{Owns} & \thead{Built in} \\
  \midrule
  \code{can_def} & Shared constants and subtypes every other block agrees on & \chapref{c10:ch} \\
  \code{meta_prev} & Getting an asynchronous signal safely into this clock domain, with no idea what
    the signal means & \chapref{c12:ch} \\
  \code{bit_timer} & \emph{When} to sample or shift, with no idea what is being sent &
    \chapref{c12:ch} \\
  \code{crc15} & The CRC-15 computation, one bit at a time, with no idea what a ``frame'' is &
    \chapref{c13:ch} \\
  \code{tx_shift_reg} / \code{rx_shift_reg} & Bit stuffing and destuffing, with no idea what an ID or
    a data byte is & \chapref{c14:ch}--\ref{c15:ch} \\
  \code{can_controller} & The only block that knows the actual CAN frame format, arbitration included
    & \chapref{c11:ch} (ports), \chapref{c16:ch}--\ref{c17:ch} (behaviour) \\
  \code{register_bank} & Pulses to pollable levels, and the register map's semantics; no idea what SPI
    or CAN is & \chapref{c18:ch} \\
  \code{spi_slave} & Bytes off a serial wire, safely into the 50 MHz domain; no idea what a register
    is & handed out, read in \chapref{c18:ch} \\
  \code{spi_reg_bridge} & The five-byte transaction; no idea what any individual register means &
    \chapref{c19:ch} \\
  \code{can_spi_node} & The connections, and the pins the bus lives on. No logic of its own &
    \chapref{c19:ch} \\
\end{booktable}

That is the whole node: four small blocks with one responsibility each and one shared package composed
inside \code{can_controller}, and then three more layers in front of it, each knowing only the one
below it. \code{can_spi_node} stands outside them all, since adapting to a particular board is a
different matter from speaking CAN, or SPI, or registers.

Read the table top to bottom and it is also a list of what would have to change to reach a different
system. Swap \code{spi_slave} and \code{spi_reg_bridge} for an I2C pair and everything beneath them is
untouched. Swap \code{can_controller} for another protocol engine and the register map moves but the
transport does not. That is what the boundaries bought.

\paragraph{The register map in summary}
Every port on \code{can_controller}'s register side corresponds to a register the parallel class's C++
driver reads or writes; the remaining five (\code{clock}, \code{reset_n} and the three bus signals)
are physical and no driver ever sees them. The correspondence is no longer a plan, as it was when
\secref{c11:sec:regmap} first drew this table \textemdash{} \code{register_bank} is the module that
realises it, so read the right-hand column as ``the register the bank presents at that index'':

\begin{booktable}{@{}p{14em}L@{}}
  \thead{\code{can_controller} port} & \thead{Register in the map} \\
  \midrule
  \code{tx_id}, \code{tx_dlc}, \code{tx_data} & \code{TX_ID}, \code{TX_DLC},
    \code{TX_DATA_LO}/\code{HI} \\
  \code{tx_req}, \code{tx_done} & \code{TX_SEND}, one bit in \code{STATUS} \\
  \code{rx_id}, \code{rx_dlc}, \code{rx_data} & \code{RX_ID}, \code{RX_DLC},
    \code{RX_DATA_LO}/\code{HI} \\
  \code{rx_valid} & one bit in \code{STATUS}, cleared via \code{RX_ACK} \\
  \code{error} & \code{ERROR_FLAGS} \\
\end{booktable}

That correspondence is where the driver work begins: everything the parallel class writes against is in
appendix~\ref{app:protocol}, and nothing in this design's VHDL.
